\subsection{Колебание бесконечной струны}

\subsubsection*{Свободные колебания: задача (1)}

\stepcounter{equation}% уникальный якорь hyperref: \tag сам счётчик не двигает
\begin{equation}\tag{1}\label{eq:inf-free}
  \begin{cases}
    u_{tt} - a^2 u_{xx} = 0, & x \in \R,\ t > 0, \\[2pt]
    u|_{t=0} = \varphi(x), \quad u_t|_{t=0} = \psi(x).
  \end{cases}
\end{equation}

\begin{warn}
На доске в первой записи коэффициент выглядит как $v^2$, во всех последующих —
как $a^2$; принято $a^2$. Там же в условии написано «$t = \infty$» — очевидная
описка вместо $t > 0$.
\end{warn}

Решение задачи~\eqref{eq:inf-free} получается заменой через характеристики —
той же самой $\xi = x - at$, $\eta = x + at$, что подробно разобрана ниже для
задачи~\eqref{eq:inf-forced}:
\[
  u(x,t) = \frac12\Big[\varphi(x-at) + \varphi(x+at)\Big]
         + \frac{1}{2a}\int\limits_{x-at}^{x+at} \psi(y)\, dy
\]
\begin{center}
\textbf{— формула Даламбера—Эйлера.}
\end{center}

\begin{warn}
На доске пределы интеграла от $\psi$ в этом месте записаны как $-\infty$ и
$+\infty$. Из итоговой формулы Даламбера (см.~\eqref{eq:dalembert}) видно, что
пределы — $x-at$ и $x+at$.
\end{warn}

\begin{idea}
Откуда берётся эта формула: оператор факторизуется,
\[
  \partial_{tt} - a^2 \partial_{xx}
  = (\partial_t - a\partial_x)(\partial_t + a\partial_x),
\]
поэтому общее решение однородного уравнения — сумма двух бегущих волн
$u = F(x-at) + G(x+at)$. Подстановка начальных условий даёт
\[
  F + G = \varphi, \qquad a\,(G' - F') = \psi,
\]
откуда и получаются полусумма сдвигов и интеграл от $\psi$.
\end{idea}

\begin{example}[начальное отклонение, нулевая скорость]
$\psi \equiv 0$, $\varphi$ — «треугольник». Тогда
$u(x,t) = \tfrac12[\varphi(x-at) + \varphi(x+at)]$: начальный горб распадается на
две волны половинной высоты, уходящие налево и направо со скоростью $a$.
\end{example}

\inputfig{figures/fig-wave-split}

\begin{example}[нулевое отклонение, точечный удар]
$\varphi \equiv 0$, $\psi(x) = \delta(x)$. Тогда
\[
  u(x,t) = \frac{1}{2a}\int\limits_{x-at}^{x+at} \delta(y)\, dy
  = \begin{cases}
      \dfrac{1}{2a}, & |x| < at, \\[6pt]
      0, & |x| > at,
    \end{cases}
\]
то есть «ступенька» высоты $\tfrac{1}{2a}$, края которой — точки $-at$ и $at$ —
разбегаются со скоростью $a$. Разрывный профиль — это, разумеется, просто
математическая модель: реальная струна таких углов не имеет.
\end{example}

\inputfig{figures/fig-delta-box}

\begin{idea}
Разница с уравнением теплопроводности видна сразу: там точечный источник мгновенно
«размазывается» по всей прямой (рис.~\ref{fig:heat-kernel}), здесь возмущение строго
заперто в интервале $|x| < at$. Это и есть конечность скорости распространения.
\end{idea}

\subsubsection*{Вынужденные колебания при нулевых данных: задача (2)}

\stepcounter{equation}
\begin{equation}\tag{2}\label{eq:inf-forced}
  \begin{cases}
    u_{tt} - a^2 u_{xx} = f(x,t), & x \in \R,\ t > 0, \\[2pt]
    u|_{t=0} = 0, \quad u_t|_{t=0} = 0 .
  \end{cases}
\end{equation}

Заметим сразу: $u|_{t=0} = 0 \Longrightarrow u_x|_{t=0} = 0$ (дифференцируем
тождественный нуль вдоль прямой $t=0$).

\paragraph{Характеристические переменные.} Положим
\[
  \begin{cases} \xi = x - at, \\ \eta = x + at, \end{cases}
  \qquad\Longleftrightarrow\qquad
  \begin{cases} x = \tfrac12(\xi + \eta), \\[2pt] t = \tfrac{1}{2a}(\eta - \xi). \end{cases}
\]
В этих переменных
\[
  u_{tt} - a^2 u_{xx} = -4a^2 u_{\xi\eta},
  \qquad\text{так что}\qquad
  -4a^2 u_{\xi\eta} = \widetilde f(\xi,\eta),
\]
где $\widetilde f(\xi,\eta) = f\!\left(\tfrac12(\xi+\eta),\ \tfrac{1}{2a}(\eta-\xi)\right)$.

Прямая $t = 0$ переходит в прямую $\xi = \eta$. Начальные условия дают на ней
\[
  u|_{\xi=\eta} = 0, \qquad u_\xi|_{\xi=\eta} = 0, \qquad u_\eta|_{\xi=\eta} = 0 .
\]

\begin{idea}
Последние два равенства — это пересчёт условий $u_x|_{t=0} = 0$ и $u_t|_{t=0}=0$:
\[
  u_x = u_\xi + u_\eta, \qquad u_t = a\,(u_\eta - u_\xi),
\]
и если обе левые части равны нулю на прямой $t=0$, то там $u_\xi = u_\eta = 0$.
\end{idea}

\paragraph{Интегрирование по характеристическому треугольнику.}
Зафиксируем точку $(\xi_0, \eta_0)$ и возьмём треугольник
\[
  D = \{\, (\xi,\eta) : \ \xi_0 \le \xi \le \eta_0, \ \ \xi \le \eta \le \eta_0 \,\}.
\]

\inputfig{figures/fig-char-xieta}

Тогда
\begin{align*}
  \iint\limits_D u_{\xi\eta} \, d\xi\, d\eta
  &= \int\limits_{\xi_0}^{\eta_0} \!\! \int\limits_{\xi}^{\eta_0}
     u_{\xi\eta}\, d\eta\, d\xi
   = \int\limits_{\xi_0}^{\eta_0} u_\xi \Big|_{\eta=\xi}^{\eta=\eta_0} d\xi
   = \int\limits_{\xi_0}^{\eta_0} u_\xi(\xi, \eta_0)\, d\xi \\[4pt]
  &= u(\xi,\eta_0)\Big|_{\xi=\xi_0}^{\xi=\eta_0}
   = -\,u(\xi_0, \eta_0),
\end{align*}
поскольку $u_\xi|_{\eta=\xi} = 0$, а $u(\eta_0,\eta_0) = 0$ (эта точка лежит на
прямой $\xi=\eta$). С другой стороны $u_{\xi\eta} = -\tfrac{1}{4a^2}\widetilde f$,
значит
\[
  \boxed{\ u(\xi_0,\eta_0) = \frac{1}{4a^2} \iint\limits_D \widetilde f(\xi,\eta)\,
  d\xi\, d\eta \ }
\]

\paragraph{Возврат к переменным $(x,t)$.} Якобиан замены:
\[
  J = \det \begin{pmatrix} 1 & -a \\ 1 & a \end{pmatrix} = 2a ,
  \qquad d\xi\, d\eta = 2a \, dy \, d\tau .
\]
Поэтому
\[
  u(x_0,t_0) = \frac{1}{4a^2} \iint\limits_{D} 2a\, f(y,\tau)\, dy\, d\tau
             = \frac{1}{2a} \iint\limits_{D} f(y,\tau)\, dy\, d\tau ,
\]
где теперь $D$ — характеристический треугольник в плоскости $(x,t)$ с вершиной
$(x_0,t_0)$ и боковыми сторонами $|x - x_0| = a(t_0 - t)$:
\[
  u(x_0,t_0) = \frac{1}{2a} \int\limits_0^{t_0}
  \int\limits_{x_0 - a(t_0-\tau)}^{\,x_0 + a(t_0-\tau)} f(y,\tau)\, dy\, d\tau .
\]

\inputfig{figures/fig-char-triangle}

\begin{warn}
И на доске, и в тетради пределы внутреннего интеграла записаны короче — как
$x_0 - at$ и $x_0 + at$ (а ниже, в формуле Даламбера, как $x - a\tau$ и
$x + a\tau$). Поскольку по внешней переменной идёт интегрирование, формально
это неверно: по рисунку характеристического треугольника срез на уровне $\tau$
имеет концы $x_0 \mp a(t_0 - \tau)$, и именно так записано выше. Судя по тому,
что сокращение повторяется во всех записях, это устойчивая скоропись лектора,
а не другая формула.
\end{warn}

\subsubsection*{Общая задача (3) и формула Даламбера}

\stepcounter{equation}
\begin{equation}\tag{3}\label{eq:inf-full}
  \begin{cases}
    u_{tt} - a^2 u_{xx} = f(x,t), & x \in \R,\ t > 0, \\[2pt]
    u|_{t=0} = \varphi(x), \quad u_t|_{t=0} = \psi(x) .
  \end{cases}
\end{equation}

По линейности задача~\eqref{eq:inf-full} есть сумма задач~\eqref{eq:inf-free}
и~\eqref{eq:inf-forced}, поэтому

\stepcounter{equation}
\begin{equation}\tag{Д}\label{eq:dalembert}
  \boxed{
  \begin{aligned}
    u(x,t) &= \frac12\Big[\varphi(x-at) + \varphi(x+at)\Big]
            + \frac{1}{2a}\int\limits_{x-at}^{x+at} \psi(y)\, dy \\[4pt]
           &\quad + \frac{1}{2a} \int\limits_0^t
             \int\limits_{x-a(t-\tau)}^{\,x+a(t-\tau)} f(y,\tau)\, dy\, d\tau
  \end{aligned}}
\end{equation}

\begin{center}
\textbf{— формула Даламбера.}
\end{center}

\subsubsection*{Компактная запись}

Обозначим интервал
\[
  \Bxat{x}{at} = (x - at,\ x + at).
\]
Заметим, что
\[
  \frac12\Big[\varphi(x-at) + \varphi(x+at)\Big]
  = \frac{1}{2a}\,\dd{}{t} \int\limits_{x-at}^{x+at} \varphi(y)\, dy
  = \frac{1}{2a}\,\dd{}{t} \int\limits_{\Bxat{x}{at}} \varphi(y)\, dy ,
\]
поскольку дифференцирование интеграла по пределам даёт
$a\varphi(x+at) + a\varphi(x-at)$. Поэтому формула Даламбера записывается
единообразно:
\[
  \boxed{\ u(x,t) = \frac{1}{2a}\left[
    \dd{}{t} \int\limits_{\Bxat{x}{at}} \varphi(y)\, dy
    + \int\limits_{\Bxat{x}{at}} \psi(y)\, dy
    + \int\limits_0^t \int\limits_{\Bxat{x}{a(t-\tau)}} f(y,\tau)\, dy\, d\tau
  \right]\ }
\]

\begin{warn}
В записи интервала на доске стояло $\Bxat{x}{at} = (x+at,\, x+at)$; в тетради
записано правильно, $(x-at,\, x+at)$, — значит на доске просто описка.
А вот в третьем слагаемом и доска, и тетрадь пишут $\Bxat{x}{a\tau}$: это то же
сокращение, что и выше, по характеристическому треугольнику там
$\Bxat{x}{a(t-\tau)}$.
\end{warn}

\begin{idea}
Запись через $\Bxat{x}{at}$ не косметическая: в размерности $n = 3$ (см.~\S3)
формула будет выглядеть ровно так же, только $\Bxat{x}{at}$ станет шаром радиуса
$at$, а интегралы — средними по его границе. Одномерный интервал — это «шар» при
$n=1$.
\end{idea}

\subsubsection*{Существование классического решения}

\begin{theorem}\label{thm:dalembert-exist}
Пусть
\[
  \varphi \in C^2(\R), \qquad \psi \in C^1(\R), \qquad f,\ f_x \in C(\R) .
\]
Тогда классическое решение задачи~\eqref{eq:inf-full} существует и выражается
формулой Даламбера~\eqref{eq:dalembert}.
\end{theorem}

\begin{proof}
Уже доказано: проведённый выше вывод и есть доказательство. Остаётся заметить,
что при указанной гладкости формула Даламбера дважды непрерывно
дифференцируема, то есть даёт \emph{классическое} решение.
\end{proof}

\begin{idea}
Требования по гладкости ровно такие, чтобы формула Даламбера после двукратного
дифференцирования дала непрерывные $u_{tt}$ и $u_{xx}$.

\medskip
$\varphi \in C^2$ — входит без интегрирования, поэтому теряет обе производные
целиком.

\medskip
$\psi \in C^1$ — стоит под интегралом с пределами, зависящими от $x$ и $t$.
Дифференцирование по пределам «съедает» одну производную: первая получается
бесплатно, на вторую нужна гладкость самой $\psi$.

\medskip
$f, f_x \in C$ — $f$ стоит под двойным интегралом, поэтому по $t$ запаса
хватает; но при дифференцировании по $x$ один раз интеграл снимается
пределами, а для второго раза нужна уже непрерывность $f_x$. Отсюда и
несимметричность условия по $x$ и $t$.
\end{idea}

\subsubsection*{Область зависимости}

\inputfig{figures/fig-dependence}

\begin{warn}
В тетради треугольник подписан «область влияния \textbf{на точку} $(x,t)$» — с
этим уточнением всё согласуется: речь именно о том, что влияет \emph{на} точку,
то есть об области зависимости. На доске слова «на точку» потерялись, оттого и
возникала путаница с конусом \emph{над} точкой. Для справки, стандартная
терминология:
треугольник \emph{под} точкой $(x,t)$ — то, от чего значение $u(x,t)$ зависит, —
называется \textbf{областью зависимости}, а областью влияния точки называют конус
\emph{над} ней.
\end{warn}

\begin{idea}
Содержательно: значение $u(x,t)$ определяется данными $\varphi, \psi$ только на
интервале $\Bxat{x}{at}$ и значениями $f$ только внутри характеристического
треугольника. Изменим данные вне этого треугольника — значение $u(x,t)$ не
изменится.

\medskip
Это и есть \textbf{конечная скорость распространения}: события в точке $(x,t)$
могут не зависеть от более ранних событий, если те произошли достаточно далеко.
У уравнения теплопроводности скорость распространения бесконечна — там
возмущение из любой точки мгновенно сказывается всюду
(ср.~рис.~\ref{fig:heat-kernel}).
\end{idea}

\subsubsection*{Непрерывная зависимость от данных}

Пусть двум наборам данных отвечают два решения:
\[
  \varphi_1, \psi_1, f_1 \ \rightsquigarrow\ u_1,
  \qquad
  \varphi_2, \psi_2, f_2 \ \rightsquigarrow\ u_2 .
\]
Введём
\[
  \delta_\varphi := \sup_x \abs{\varphi_1(x) - \varphi_2(x)}, \qquad
  \delta_\psi := \sup_x \abs{\psi_1(x) - \psi_2(x)},
\]
\[
  \delta_{f,T} := \sup_{\substack{x \in \R \\ t \in [0,T]}}
    \abs{f_1(x,t) - f_2(x,t)} .
\]

\begin{theorem}
Классическое решение задачи~\eqref{eq:inf-full} непрерывно зависит от
$f, \varphi, \psi$ в том смысле, что для любого $T > 0$ выполняется неравенство
\[
  \abs{u_1(x,t) - u_2(x,t)} \;\le\; \delta_\varphi + \delta_\psi\, T
    + \delta_{f,T}\,\frac{T^2}{2}
  \qquad \forall\, x\in\R,\ t \in [0,T] .
\]
\end{theorem}

\begin{proof}
Оставлено как домашнее задание; сама выкладка разобрана в пояснении ниже.
\end{proof}

Единственность здесь не доказывается: она будет получена позже и сразу для
$n$-мерного случая (\S4).

\begin{idea}
Оценка читается прямо из формулы Даламбера: разность решений — это та же формула,
применённая к разностям данных. Полусумма сдвигов даёт $\delta_\varphi$; интеграл
от $\psi$ по интервалу длины $2at$ даёт
$\frac{1}{2a}\cdot 2at \cdot \delta_\psi = \delta_\psi\, t$; двойной интеграл по
треугольнику площади $a t^2$ даёт
$\frac{1}{2a}\cdot a t^2 \cdot \delta_{f,T} = \delta_{f,T}\,\frac{t^2}{2}$. Затем
$t$ оценивается через $T$.

\medskip
На доске эта выкладка сначала записана с $t$, а в формулировке теоремы — уже с
$T$; это одно и то же неравенство.

\medskip
Вместе с теоремой~\ref{thm:dalembert-exist} и единственностью (она следует из
самой формулы: решение задаётся данными однозначно) это даёт
\textbf{корректность по Адамару} задачи~\eqref{eq:inf-full}: существование,
единственность, непрерывная зависимость от данных.
\end{idea}
